\documentclass[11pt]{article}
\usepackage[margin=1in]{geometry}
\usepackage[T1]{fontenc}
\usepackage{lmodern}
\usepackage{amsmath,amssymb,amsthm,booktabs,tabularx,array}
\usepackage[hidelinks]{hyperref}
\usepackage{microtype}
\usepackage{enumitem}
\setlist{nosep,leftmargin=*}
\setlength{\parskip}{4pt}
\setlength{\parindent}{0pt}
\setlength{\emergencystretch}{2em}
\clubpenalty=10000
\widowpenalty=10000
\newtheorem{theorem}{Theorem}
\newcommand{\code}[1]{\texttt{\small #1}}
% Replace only with author-supplied names and affiliations before submission.
% No author identity, affiliation, or license has been inferred.
% User confirmed sole authorship on 5 November; exact publication name pending.
\author{}
\date{5 November 2026}
\newcommand{\publicationstatus}{Preparation draft: author metadata pending; not yet submitted.}

\title{TrustBound: A Mechanized Case Study\\of Payment-Epoch Completeness}
\begin{document}
\maketitle
\begin{center}\small\publicationstatus\end{center}
\begin{abstract}
A payment receipt establishes an observed outcome, not that every payment has
been accounted for. We study exact closure of a predefined payment epoch when
agents submit requests to a trusted payment service. Recording an operation
before granting execution authority and serializing both with an admission
fence yield a fixed journal covering past effects and remaining execution
rights in a specified transition system. Truthful per-identity outcomes and
irreversible discharge of those rights then establish exact closure for every
permitted finite continuation. A second model permits unknown inert provider
objects, provided every activation is mediated. We mechanize these conditional
results in Lean. Reproduction checks 14 exported results and rejects eight
premise-removal mutations. A finite evidence-rule diagnostic generates 2,048
records over 64 distinct public inputs; two labels for the same reconciliation
rule agree by construction. This is not an independent baseline implementation
or a validation of the formal model. Historical processor tests illustrate selected
behaviors but do not establish real-world admission completeness. The
contribution is a checked construction, a constrained model mapping, and a
reproducible account of the completeness boundary. We do not establish a new
general closure principle, weaker trust than the closest work, or a verified
payment deployment.
\end{abstract}

\section{Introduction}
A travel business may delegate flight and hotel booking to separate agents.
A flight payment can be known while a hotel payment has an uncertain response.
The service may stop admitting new work while the earlier hotel request remains
able to execute. Booking completion, a balanced reservation ledger, and a signed
list of observed payments each leave a question unanswered: is the relevant
payment history complete and unable to grow under earlier authority?

This paper fixes one merchant payment path and epoch. Agents are request
proposers; the service, its journal, and the relevant source contracts are
trusted. Travel is a motivating application, not a claim of verified composition
across booking providers or independent issuers. Removing an itinerary
coordinator does not remove trust in the payment authority.

We make three contributions within that scope. First, a guarded admission
model derives fixed-population coverage and effect stability. Checked
counterexecutions distinguish an admission fence from drainage and show the
effect of an external bypass. Second, a provider-contract model demonstrates
that unregistered objects need not be enumerated if they are inert and every
effect-capable activation is mediated. Third, an evidence-only finite diagnostic
separates its prescribed certificate predicate from bypass violations and
exposes the limits of authentic but incomplete evidence. Its conventional arm
is the same acceptance function under another label, not a separate system.

These are claims about this formalization and case study. Fencing, conservative
recovery, budget reservations, and stable-property reasoning are not claimed as
new. Section~\ref{sec:related} compares the closest work under comparable trust
assumptions. The experiment establishes no superiority over reconciliation.

\section{Scope and observation contract}
\label{sec:scope}
\subsection{Trusted-service assumptions}
Agents may replay, reorder, forge, or omit requests, but cannot access the
service machine, payment keys, journal, or privileged controls. The prototype's
request interface accepts previously approved command identities and limited
status/reconciliation requests. General approval and multi-tenant authorization
are outside this evaluation.

The epoch's economic scope must be fixed before execution. A second credential,
operator, or automation that can create an attributable effect is an in-scope
route to mediate, not a reason to redefine the scope after observing an effect.
Recovery of an admitted request remains part of its execution authority.
The modeled epoch starts empty. Applying the theorem to a new epoch therefore
requires that no earlier right or effect attributable to that epoch is omitted
from the initial abstraction. Resetting a journal or relabeling a delayed
request does not meet this requirement. Importing pre-existing work would
require an explicitly justified initial-state invariant, which is not provided
by the empty-start theorem here.

\begin{description}[style=nextline]
\item[A1: Request-only agents.] Agents have no access to the trusted service's
internals or administration. Same-host arbitrary-code adversaries are excluded.
\item[A2: Exclusive scoped activation.] Every effect-capable route in the
predefined epoch is mediated. Identities cannot be reassigned to evade the fence.
\item[A3: Durable ordered service.] Recording precedes authority, recording
and grants serialize with the fence, and committed state cannot roll back.
\item[A4: Truthful source terminality.] Outcomes bind the correct identities
and facts. Discharge irreversibly removes all remaining effective authority,
including in-flight delivery and provider recovery.
\item[A5: Faithful complete export.] The auditor receives the authoritative
journal and truthful outcomes for the same declared scope and cut.
\end{description}

A1 does not imply A2. The state machines enforce abstract ordering in A3; no
refinement from executable Python and storage to those machines is proved.
A4 is a semantic source contract, not a conclusion from a terminal-looking
client status. A signature authenticates an assertion without proving A5.
A separately obtained view from the same processor is not an independently
trusted issuer. Successful API tests do not establish these premises universally.

\subsection{Four different assertions}
\emph{Admission closed} means no new records or grants are accepted. Earlier
rights may still execute. \emph{Exposure covered} means retained liability
covers possible cumulative debits; an uncertain operation can meet a bound
without a known outcome. \emph{Execution stopped} means no permitted
continuation adds an attributable effect, although past effects may remain
unknown to the observer. \emph{Exactly closed} combines complete accounting
of the admitted population with exclusion of future effects.

Our main theorem concerns effect identities. It does not establish booking
success, legal finality, eventual completion, or a net balance immune to refunds
and disputes. Separate amount and dispatch models in the research workspace
are not composed into an end-to-end monetary theorem here.

\section{Admission, fencing, and exact closure}
\label{sec:theorem}
\subsection{State and transitions}
The Lean state is $(J,I,R,E,F)$: the journal, issued identities, remaining
effective rights, effect history, and closed flag. The first four components
are lists of natural-number identities; the initial lists are empty and $F$ is
false. Set notation below abbreviates list membership, except that history
equality is actual list equality. $R$ includes external delivery authority,
not merely the client's queue. $E$ is ghost history, not an auditor input.

\begin{table}[ht]
\centering\small
\begin{tabularx}{\linewidth}{@{}l X X@{}}
\toprule
Action & Guard & State change \\
\midrule
Record $i$ & $\neg F$, $i\notin J$ & Prepend $i$ to $J$. \\
Grant $i$ & $\neg F$, $i\in J$, $i\notin I$ & Prepend $i$ to $I$ and $R$. \\
Deliver $i$ & $i\in R$ & Remove $i$ from $R$; prepend it to $E$. \\
Cancel $i$ & No additional guard & Remove $i$ from $R$. \\
Fence & No additional guard & Set $F$ to true. \\
\bottomrule
\end{tabularx}
\caption{The admission model. Cancel removes current effective rights; it is
not an API cancellation acknowledgement or a permanent identity tombstone.
Unlisted fields are unchanged.}
\label{tab:steps}
\end{table}

All guards and updates in Table~\ref{tab:steps} are atomic abstract transitions.
The issued list prevents granting the same identity again. The model represents
one effect per identity, not arbitrary partial-capture or recovery lineages.
A finite run is the reflexive transitive closure of the transition relation.
The proofs place no finite bound on run length or on the identity domain.

Cancel does not add to $I$ or set a terminal flag. Recording $i$, cancelling it
before any grant, and then granting and delivering it is permitted while the
epoch is open. After a grant, $I$ prevents regrant; after a fence, all grants
are disabled. Thus a cancelled identity is not unconditionally terminal in
this model. Permanent terminal-identity exclusion is an additional rule of
the provider model in Section~\ref{sec:mapping}. Neither cancellation nor
discharge erases a past effect or represents a refund.

\subsection{Conditional theorem}
\begin{theorem}[Fenced identity coverage and exact closure]
Let $s$ be reachable from the initial state with $F(s)=\mathrm{true}$, and let
$t$ be reachable from $s$ by any finite permitted run. Then
\begin{align}
J(t)&=J(s),\label{eq:fixed}\\
R(t)\cup E(t)&\subseteq J(s).\label{eq:covered}
\end{align}
If every $i\in J(s)$ is absent from $R(s)$ and a predicate $\mathrm{reported}$
satisfies $\mathrm{reported}(i)\Leftrightarrow i\in E(s)$ for all such $i$, then
\begin{align}
R(t)&=\varnothing,& E(t)&=E(s),\label{eq:stable}\\
i\in E(t)&\ \Longleftrightarrow\ i\in J(s)\land\mathrm{reported}(i).
\label{eq:exact}
\end{align}
\end{theorem}

\begin{proof}
Every transition preserves coverage $R\cup E\subseteq J$: a grant needs a
journal entry, delivery consumes an existing right, and cancellation introduces
no effect. Reachability from the empty state therefore establishes coverage.
After fencing, record and grant are disabled, so the journal is fixed. This
proves~\eqref{eq:fixed}--\eqref{eq:covered}. Coverage and discharge of every
journaled identity imply $R(s)=\varnothing$. No subsequent transition can
reintroduce rights or deliver an effect, yielding~\eqref{eq:stable}. Truthful
reports identify precisely the executed journal entries, giving~\eqref{eq:exact}.
\end{proof}

The corresponding checked declarations are
\code{fenced\_population\_immutable}, \code{post\_fence\_coverage},
\code{drained\_fence\_effects\_stable}, and
\code{reconciled\_fence\_complete} in \code{Admission.lean}.
Future stability needs drainage but not an outcome report; exact identification
of the past additionally needs truthful reports. The theorem derives coverage
from guarded transitions. It does not discover whether every real route obeys
those transitions or whether a supplied export is faithful.
Truthful Boolean reports determine effect membership, not the historical
ordering of effects. The list equality states that the modeled history remains
unchanged; it does not reconstruct that ordering from the reports.

\subsection{Counterexecutions and transfer obligation}
The checked \code{fence\_without\_drain} execution records and grants an
identity, closes admission, and then delivers it. With a completed flight and
pending hotel operation, for example, $J=\{f,h\}$, $E=\{f\}$ and $R=\{h\}$.
Authorization may be safe while exact closure is false. This is not an attack
on the guarantees of the closest papers.

The checked \code{bypass\_breaks\_coverage} example adds an effect outside the
transition relation without changing the journal. This violates A2, rather
than disproving the theorem under its premises. Richer provider queries can
distinguish such histories; the example makes no full-API indistinguishability
claim.

The generic \code{refinement\_transfers\_coverage} result states the remaining
transfer obligation explicitly. Every real action must simulate zero or more
model steps; every real effect and right must be represented; the initial
abstraction must be initial; and the exported journal must equal the abstract
journal. We do not prove those premises for an executable deployment.
The bridge transfers membership coverage only. Its representation premises
are one-way implications and do not require a one-to-one correspondence
between concrete effect occurrences and abstract identities. Consequently it
does not transfer exact event counts, ordering, or monetary totals, and is not
an end-to-end refinement proof for Theorem~1. Such a result would additionally
need an outcome-preserving representation and a proof that concrete drainage
and observation facts imply the theorem's premises.

\section{Constrained provider-contract mapping}
\label{sec:mapping}
The second model, \code{StripeGate.lean}, adds a collection of provider objects
and terminal identities around the admission state. Allocation creates an
inert object. Registration requires an existing object and an open epoch.
Dispatch requires registration, an open epoch, an unused identity, and absence
of terminal status. A paid transition consumes a right, records an effect,
and marks terminality. Cancellation removes the right and marks terminality.
Lost replies, reads, and rejected attempts can leave the model state unchanged;
any still-effective authority must remain represented in $R$.

\code{provider\_step\_refines} and \code{provider\_runs\_refine} map these
specified transitions to the admission model. Reachability then implies that
an unregistered object has neither a right nor an effect.
\code{terminal\_manifest\_drains} establishes an empty rights list when all
journal entries are terminal. After a fence, \code{terminal\_cut\_stable}
gives history stability. Terminality here is defined by permanent exclusion
of future effects, not inferred from a status name.

The practical motivation is an unconfirmed object with manual confirmation
and capture, durably registered before activation. Unknown inert objects need
not be enumerated to cover payments. The hard obligation is showing that
another credential, operator, automation, or configuration cannot activate one.
A lost creation response is harmless only under this inert-creation contract.
This is a model-to-model mapping, not verification of Stripe's servers, an
adapter, or a weaker-trust alternative to complete effect mediation.
Allocation can remain possible after a fence without invalidating the result:
registration and activation for that epoch are blocked. Conversely, a single
concrete operation that can debit more than once or create separately effective
descendants cannot be mapped to one modeled identity merely by sharing a
request identifier. Its effect granularity and remaining authority need a
separate justified representation; this paper supplies no such mapping.

\section{Evaluation and reproduction}
\label{sec:eval}
We separate fresh checks of the selected core from historical observations.
On 5 November 2026, a source-isolated run copied eight selected source and
configuration files, using existing Lean 4.19.0, Python 3.12.14, and pytest 9.0.2
installations. It made no provider requests. The release artifact was also
extracted into a separate directory and its included runner was executed there
with the same installed tools. Neither check was an independent clean-machine
installation or a real-deployment verification.

\subsection{Formal checks}
All eight exported admission results and all six provider-contract results
passed. Their accepted axiom reports use only \code{propext} and
\code{Quot.sound}, or no axioms; no custom axiom or unfinished proof is accepted.
The provider checker recompiles its admission dependency, which is not counted
as eight additional results.
The 14 exports include invariant and mapping results, two counterexecution
statements, and a conditional transfer bridge; they are not 14 independent
security guarantees. The two Lean modules use sequential atomic abstract
steps. Arbitrary interleavings of those steps are represented, but concrete
concurrency, storage linearizability, and crash recovery are not thereby proved.

The checkers separately remove eight premises or effects: record-before-grant,
record/fence ordering, grant/fence ordering, delivery authorization, inert
creation, registered dispatch, terminal dispatch exclusion, and source discharge
on cancellation. All eight altered proofs are rejected. This shows dependence
of the existing proof text on the alterations, not universal premise minimality.
Expected rejection logs are distinct from accepted proof logs. They may contain
Lean's error-recovery markers and are not counted as accepted proofs.

\subsection{Equal-information evidence-rule diagnostic}
The generator fixes two labels, flight and hotel, each with one admission.
For each label it selects one of four tuples $(\text{public status},
\text{liability},\text{hidden debit})$:
$(\mathrm{pending},100,0)$, $(\mathrm{pending},100,75)$,
$(\mathrm{confirmed},75,75)$, or $(\mathrm{released},0,0)$.
It independently varies two report-presence bits, two fence bits, and one
bypass bit. This gives $4^2\cdot2^2\cdot2^2\cdot2=512$ symbolic cases per arm,
half with mediation and half with an unreported in-scope bypass effect.
Four arms yield 2,048 records. Authentication is idealized.
There are only 64 distinct public evidence inputs per arm: each path has
four report possibilities (omitted, pending, confirmed, released) and two
fence possibilities. Hidden alternatives and arm labels reuse these inputs.
The records are Cartesian assignments, not 2,048 independent executions or
reachable traces of the Lean transition system.

All arms receive the same public evidence: path and epoch identifiers, claimed
fences, high-water marks, and reported rows. The mediation label reaches only
the grader; hidden debit amounts are recorded but not used in either verdict
or grading. Each arm checks path/epoch and row validity;
an exact-closure verdict also rejects a pending row. The two strong arms require
a fence and a contiguous row sequence through each high-water mark.
Signed-closure and witnessed-log arms omit those latter checks. These two weak
arms are mechanism ablations, not implementations of deployed products or of
the cited papers.
The two weak labels also share a function; no separate signature verification
or witnessing protocol is implemented. Their counts illustrate omitted checks,
not a demonstrated weakness of signatures or transparency logs as systems.

For grading, coverage is true exactly when no bypass occurs, both reports are
present, and both fences hold. Exact closure additionally requires both statuses
to be non-pending. This is the finite model's target predicate, not a measured
account-wide property or an independently verified monetary condition.
In particular, this grading predicate is not the Lean membership
invariant $R\cup E\subseteq J$: the latter can hold in an open epoch or with
a missing export row. The diagnostic's field named \code{coverage\_true}
means complete fenced evidence by stipulation, not that invariant. Hidden
debit amounts are saved for illustration but are not inputs to either grading
predicate. Thus the zero-mismatch mediated rows test consistency with the
prescribed evidence rule, not independent semantic soundness of Theorem~1.
No refinement from this two-path diagnostic to the one-epoch Lean model is
proved, and no composition of independent payment authorities is evaluated.
The saved field \code{assumptions\_hold} records only the absence of a bypass;
it does not assert all A1--A5. Rows in that partition can still have missing
reports or open fences. Table~\ref{tab:baseline} therefore labels it
``Mediation'' rather than ``all assumptions hold.''

\begin{table}[ht]
\centering\small
\begin{tabular}{@{}llrrrr@{}}
\toprule
Arm & Premises & Cases & Accept & Mismatch & Cov./mismatch \\
\midrule
Signed closure & Mediation & 256 & 144 & 140 & 256/240 \\
Witnessed log & Mediation & 256 & 144 & 140 & 256/240 \\
Fenced reconciliation & Mediation & 256 & 4 & 0 & 16/0 \\
Conventional reconciliation & Mediation & 256 & 4 & 0 & 16/0 \\
Signed closure & Bypass & 256 & 144 & 144 & 256/256 \\
Witnessed log & Bypass & 256 & 144 & 144 & 256/256 \\
Fenced reconciliation & Bypass & 256 & 4 & 4 & 16/16 \\
Conventional reconciliation & Bypass & 256 & 4 & 4 & 16/16 \\
\bottomrule
\end{tabular}
\caption{Accepts and accepted cases violating the diagnostic's stipulated
target predicate (mismatches). The last column uses its evidence-coverage
predicate. These are not semantic soundness or failure-rate measurements.
Bypass cases violate A2.}
\label{tab:baseline}
\end{table}

The two strong arms share their acceptance rule and tie on every case, not only
in aggregate. Thus the tie is a consistency control, not an independently
implemented product comparison. A separate recount of the saved records
checks the summaries and all 512 groups' identical public evidence. The
generated result matches the corrected 1 November artifact byte for byte.
All nine existing baseline tests pass, covering evidence restrictions,
uncertainty, malformed or missing records, and authentic-but-false cuts.
This fairness control removes an information advantage but supplies no
comparative performance or security result. An independent implementation
would be needed to evaluate differences between reconciliation algorithms;
the current contribution does not claim such a difference.

An earlier 27 October comparison leaked a hidden-effect flag into the strong
verifier. Its zero-false-acceptance result is withdrawn and excluded from the
positive evidence. The corrected comparison exposes rather than detects a
bypass: identical supplied evidence cannot disclose a missing effect merely
because it carries an authentic assertion of completeness.

\subsection{Historical observations, not fresh provider validation}
The saved 3 November local simulator report records 40 assertions. Two
executions have identical known-object transcripts and gate snapshots but
processor totals of 200 and 300 cents, with an alternate activation route in
the latter. A broader simulated ledger distinguishes them. This demonstrates
a restricted-view boundary under an A2 violation, not indistinguishability
under all Stripe queries.

The saved controlled-merchant summary records 21 assertions over nine real
Stripe test-mode HTTP requests, capturing 200 test cents. Nine forged service
requests were rejected without a processor call. Previously approved work
executed after fencing; a successful capture reply was discarded. Database
reopen preserved uncertainty and rejected duplicate dispatch; a fresh read
resolved the known operation. Separately obtained PaymentIntent and Charge
lists reconciled it, using the same processor/key and non-atomic reads.
The summary explicitly reports that complete admission coverage and credential
isolation were not established. A same-user process could open the credential
and journal, failing the stronger host-isolation boundary. A1 excludes that
threat; it does not convert the failed check into success.

These historical experiments were not rerun for this paper package. Only
non-sensitive saved summaries are included, without raw private provider
records or provider replay scripts. Their observations are supporting reports,
not independently reproducible API evidence in the minimal artifact. The core
proofs and finite baseline require no provider access. No production payment
or standards-conformant end-to-end UCP/AP2 integration is claimed.

\section{Closest work and contribution boundary}
\label{sec:related}
Distributed snapshots already connect global-state observation to stable
properties~\cite{chandy}. This work does not introduce a general stable-property
or termination-detection principle; it specializes the obligations to a small
payment identity model.

\paragraph{AID-Guard.}
AID-Guard places authorization at a trusted effect boundary and conditions
recovery on provider-qualified terminal evidence and delivery fencing;
uncertainty remains reserved when those contracts are absent~\cite{aid}.
We explicitly mechanize a fenced epoch's identity population and discharge in
Lean. We do not remove AID-Guard's mediation or terminality premises or improve
its guarantee under matched assumptions. Its atomic effect scope also cannot
be treated as a promise about arbitrary provider-created partial effects.

\paragraph{Op.}
Op separates observed payment from remaining execution rights and conditions
stable reconciliation on complete local/provider evidence and absence of later
relevant changes~\cite{op}. Its closure witness binds executed and remaining
quantities. Our guarded model derives population coverage and stability, but
does not establish Op's external completeness premise for a real processor or
address its broader monetary obligations and corrections. The text inspected
on 5 November is a live author manuscript; no earlier cutoff priority is inferred.

\paragraph{HoldSpec.}
HoldSpec checks authorization-hold lifecycle properties under provider profiles
in a bounded model; its stated scope excludes concurrent-operation races, and
its reported tests do not exercise real providers~\cite{hold}. Our population
of identities and one historical test-mode path differ in scope and evidence.
Neither difference demonstrates a new terminality principle, a concurrency
solution, or comparative superiority.

These are source-level comparisons, not reproduced experiments or checked
reductions between the models, and not an exhaustive systematic literature
review. No priority or absence-of-prior-work claim follows. A permanent stop signal without the past outcome
would support future stability, already provided by our drain theorem, but
would not supply exact past accounting. A complete descendant manifest would
likewise assume, rather than resolve, a provider-created-identity coverage
obligation. Neither observation is advanced as a new security result.

\section{Limitations and conclusion}
The theorem is conditional on a correct transition abstraction. Every
effect-capable route and source-side right must be represented, including
delayed deliveries and recovery. Matching reports, signatures, known-object
reads, and successful tests do not independently discharge these obligations.
The source and the payment service remain trusted.

No end-to-end monetary composition, multi-worker crash-recovery proof,
independent-issuer completeness certificate, account-wide audit, hostile-host
security property, or finite completion deadline is established. Unresolved
operations may remain pending indefinitely. The finite baseline is neither
a throughput study nor evidence of deployment error rates.

The supported claim is that record-before-authority and a serialized fence
derive a fixed covered identity population under the specified rules; complete
source discharge and truthful reconciliation then establish exact closure for
every permitted finite continuation. The checked construction, constrained
mapping, and explicitly limited evaluation support a mechanized case study.
They do not establish a new general completeness principle or a strictly
weaker trust contract than the closest work.

\appendix
\section{Artifact and reproduction}
The accompanying minimal artifact contains the two Lean modules, their
checkers, the finite comparator and its nine tests, a reproduction runner,
reference output, and saved evidence. No payment credentials or provider SDK
are required. Use Lean 4.19.0 and Python with pytest 9.0.2; the executed Python
version was 3.12.14. From the artifact root, run:
\begin{verbatim}
python3 scripts/reproduce_core_claim.py \
  --lean /path/to/lean-4.19.0/bin/lean \
  --output results/new-run
\end{verbatim}
The output directory must not already exist. The runner copies the selected
inputs, supplies a small environment without payment credentials, disables
automatic pytest plugin loading, and writes command logs, source hashes, JUnit
output and a result summary. Expected totals are 14 exported results, eight
rejected mutations, nine passing tests, 2,048 records, and per-case strong-arm
agreement. No other historical module or full regression suite is silently
included in these counts.

Saved public logs replace local paths with declared tokens; the original
workspace evidence remains unchanged. A provenance file identifies original
and exported hashes. Hashes detect changes relative to the supplied manifest;
they are not independent witnesses of completeness or authorship. A clean
directory extraction check exercises the packaged files using the existing
toolchains, not an independent environment. Historical provider summaries are
clearly separated from the reproducible core.
Two additional reviewer examples in \code{review/AdmissionSemantics.lean}
check cancellation before a first grant and membership coverage without a
fence. They were separately compiled while reviewing the manuscript, are not
part of the 14-export count, and establish no new general security claim.

\begin{thebibliography}{99}
\bibitem{chandy}
K. Mani Chandy and Leslie Lamport.
Distributed Snapshots: Determining Global States of Distributed Systems.
\emph{ACM Transactions on Computer Systems}, 3(1):63--75, 1985.
\url{https://lamport.azurewebsites.net/pubs/chandy.pdf}.

\bibitem{aid}
Yingzhe Tong, Leyu Dai, and Songhui Guo.
AID-Guard: Stateful Authorization for Delegated Agent Effects.
arXiv:2608.21159v1, 2026.
Sections II--IV, especially III-E.
\url{https://arxiv.org/html/2608.21159v1}.

\bibitem{op}
Raeez Lorgat. Op: Compliance-Carrying Operations.
Unversioned author manuscript, September 2026; accessed 6 November 2026.
The live text may change.
Sections 5.2, 5.8--5.10, especially Proposition 5.7.
\url{https://research.momentum.inc/papers/op.html}.

\bibitem{hold}
Abhishek Sharma.
HoldSpec: A Machine-Checked Model and Conformance Suite for the Payment
Authorization Hold Lifecycle.
arXiv:2610.14419v1, 2026. Sections III-D, V, VIII.
\url{https://arxiv.org/html/2610.14419v1}.
\bibitem{quiescence}
Genliang Zhu and Chu Wang.
Authorization Revocation for Long-Running AI Agents: Root-Scoped Quiescence
under Delegation and Asynchronous Execution.
arXiv:2609.21284v1, 2026. Sections 2.4 and 5.
\url{https://arxiv.org/html/2609.21284v1}.

\bibitem{budget}
Genliang Zhu and Chu Wang.
Fault-Tolerant Budget Conservation in Distributed Multi-Agent Delegation.
arXiv:2610.00349v1, 2026. Sections 4.7--4.8 and 7.6.
\url{https://arxiv.org/html/2610.00349v1}.
\bibitem{effectbound}
Igor Santos-Grueiro.
When Does Authorization End? Effect Closure at Provider Boundaries.
arXiv:2609.02866v1, 2026. Sections II--III, V-C and Appendix C.
\url{https://arxiv.org/html/2609.02866v1}.

\bibitem{salas}
Jesus Salas.
Authority at Commit Time: Reject-and-Rerun Semantics for Governed Agentic Systems.
arXiv:2609.31490v1, 2026. Sections 2.2, 3.2, 6.2 and 7.
\url{https://arxiv.org/html/2609.31490v1}.

\bibitem{acquirebound}
Genliang Zhu and Chu Wang.
Runtime Authorization for Resources Acquired by AI Agents (AcquireBound).
arXiv:2609.14744v2, 2026. Sections 4.5, 5.6--5.7; Theorem 14.
\url{https://arxiv.org/html/2609.14744v2}.

\bibitem{protocolanalysis}
Ke Jiang, Mohan Yu, Yuan Chang, Mohit Kumar Jangid, Jianyu Niu,
Cong Wang, and Yinqian Zhang.
A Formal Analysis of Agent Payment Protocols.
arXiv:2609.00060v1, 2026. Table II and Section IV.
\url{https://arxiv.org/html/2609.00060v1}.

\bibitem{cta}
Igor Santos-Grueiro.
Temporary Authority, Permanent Effects: Commit-Time Authorization for LLM Agents.
arXiv:2607.10487v2, 2026. Sections 3--4, especially 4.2--4.3.
\url{https://arxiv.org/pdf/2607.10487v2}.

\bibitem{closuredraft}
D. O. Watts.
Revocation Closure and Authority Epochs for Agent Delegation.
Individual Internet-Draft, draft-watts-oauth-agent-revocation-closure-03,
30 October 2026. Sections 9--11; version, document date and sections verified
on Datatracker, 6 November 2026. Work in progress; not an IETF standard.
\url{https://datatracker.ietf.org/doc/html/draft-watts-oauth-agent-revocation-closure-03}.
\end{thebibliography}

\end{document}
